\cotefile{157-5}
\scanpage{85}
\mspage{72}

alors l'axiome de factorisation \inserted{$\mathrm{Cof}_3$} est satisfait. Inversement, si $\mathrm{Cof}_3$ est satisfait, et si $W$ est stable par facteurs directs, alors on a 4.4. Donc si $W$ est stable par facteurs directs, $\mathrm{Cof}_3$ équivaut à (4.4.)

4.2. Il me semble que ces conditions, en pratique, ne seront guère jamais vérifiées, \unsure{que|car} si $W$ (qui est donné en \deleted{tout} premier lieu,\inserted{\unsure{satisfaisant à} $\underline{\mathrm{Cof}_0^{*}}$}) \margin{$*$ i.e. \unsure{saturation|stabilité} faible}
% slanted
n'est \deleted{\unsure{satisfai}} pas stable par $\unsure{\varinjlim|\underline{\lim}}$ filtrantes.

À ce moment, \unsure{pour} construire la structure à cofibrations, on part de $C_0$ \unsure{petite|petites} partie de $\operatorname{Fl}M$, \unsure{disons|choisie} telle que l'on ait
\[
(C_0)_{*}\subset W,\tag*{(4.5.)}
\]
on \unsure{prend|peut} pour \inserted{au choix} $C$ le saturé \inserted{$C'_0$} de $C_0$ par cochg\textsuperscript{t}. de base et composition transfinie,, \deleted{\unsure{ou}} $C=\widetilde C_0={}$ ens des facteurs directs de $C'_0$. Les conditions $\mathrm{Cof}_1$,\inserted{(stabilité de $C$)} $\mathrm{Cof}_3$\inserted{(factorisation)} sont satisfaites automatiquement, et il reste à s'occuper de $\mathrm{Cof}_2$, qui revient à la stabilité de $C\cap W$ par cochg\textsuperscript{t}. de base (la stabilité par comp., et que ça contienne les iso, étant \unsure{immédiat|immédiate}). C'est une condition délicate, et la seule façon \inserted{\unsure{standard|standart}} que je vois pour l'assurer, c'est en supposant
\[
C_0\subset\mathrm{Cof}_W\,;\tag*{(4.6.)}
\]
% archive sheet number: 84
\scanpage{86}
\mspage{73}

la stabilité de $W$ par $\unsure{\varinjlim|\underline{\lim}}$ filtrantes (donc aussi par facteurs directs) assure alors qu'on a également
\[
C\subset\mathrm{Cof}_W,\tag*{(4.7)}
\]
et par suite
\[
\left\{\begin{aligned} C\cap W&\subset\mathrm{Cof}_W\cap W=W^{\mathrm{univ}}\\ \text{d'où}\quad \mathrm{TC}&\overset{\mathrm{déf}}{=}C\cap W=C\cap W^{\mathrm{univ}},
\end{aligned}\right.\tag*{(4.8)}
\]
ce qui assure la stabilité de $C\cap W$ par cochg. de base. La stabilité \inserted{de $\mathrm{TC}$} par comp. transfinie résulte de celle de $C$, et de la stabilité de \deleted{\unsure{$C$}}$W$ par comp. transfinie.

4.3. Je reviens aux conditions \unsure{primitives|générales} \textcircled{a} \unsure{à|et} \textcircled{\unsure{f|g}} de 4.1, sans supposer néc. $W$ stable par $\varinjlim$ filtrantes. Je veux à présent \unsure{étendre|étudier} la structure à cofibrations \unsure{aux|sur} $M(I)$, $I$ quelc. de $\mathrm{Cat}$. Nous \unsure{n'avons|n'aurons} pas pour ceci \inserted{(je crois)} à utiliser les développements des \S\S\,2,3, mais nous ferons le lien avec ceux-ci par la suite. \tikz[x=1ex,y=1ex]{\draw (0,0)--(.9,0)--(1.26,1.35); \draw (.8,.08)--(1.15,1.4); \draw (1.35,1.43)--(2.2,1.6);}

Désignons par $I_0$ la catégorie discrète ayant mêmes objets que $I$, d'où un foncteur canonique
\[
\unsure{u|v}=u_I:I_0\longrightarrow I.\tag*{(4.9)}
\]
L'idée est de définir la structure de $M(I)$ via la structure évidente de $M(\unsure{I_0})\simeq M\deleted{\unsure{\ell}}^{I_0}$, en jouant sur le couple de foncteurs adjoints
\diagram{\begin{tikzcd}[ampersand replacement=\&,column sep=huge] M(I_0) \arrow[r,shift left,"\varphi=u_!"] \& M(I) \arrow[l,shift left,"\psi=u^*"]
\end{tikzcd}\,.}
% archive sheet number: 85
\scanpage{87}
\mspage{74}

Nous allons présenter les choses de façon un peu plus générale.

\uline{Définition} 4.4. J'appelle structure de \unsure{cofibrations|cofibration} \uline{spéciale} une structure $(M,W,C)$ satisfaisant les conditions a) à f) de 4.1., \deleted{\unsure{que}} je \unsure{reprends|reformule}

\textcircled{a} $M$ stable par petites limites inductives, et les objets de $M$ sont accessibles,

\textcircled{b} $W$ est \unsure{fbt|ftt|fort} saturée, (\deleted{\unsure{de}} $\mathrm{Cof}_0$ de p. 39), $C$ stable par cochgt. de base et par compositions \deleted{\uline{transfinies},} $\mathrm{TC}\overset{\mathrm{déf}}{=}C\cap W$ stable par cochgt. de base (et par composition automatiquement)

\textcircled{b$'$} $C$ et $\mathrm{TC}$ sont stables $\unsure{=}$ par composition \uline{transfinie}

\textcircled{c} $\exists$ \uline{petite} partie $C_0\subset C$ telle que $C\subset\widetilde C_0$ ($\Longleftrightarrow\widetilde C=\widetilde C_0$)

\textcircled{d} $\unsure{\deleted{C}|{}}\mathrm{TF}\overset{\mathrm{déf}}{=}C_*\subset W$ (ce qui implique, \inserted{joint à \textcircled{a}--\textcircled{c},} l'axiome de factorisation $\mathrm{Cof}_3$ de page 39).

4.5. Nous allons donner une recette générale d'extension d'une telle structure, quand on dispose d'un couple de foncteurs adjoints
\diagram{\begin{tikzcd}[ampersand replacement=\&,column sep=huge] M \arrow[r,shift left,"\unsure{u|v}_!"] \& M' \arrow[l,shift left,"u^*"]
\end{tikzcd}\,,}
$M$ étant muni d'une structure de cofibration spéciale \inserted{$(W,C)$}, et $M'$ satisfaisant la condition liminaire \textcircled{a}. On \unsure{définira|définit} alors les parties $W',\deleted{\unsure{(C)}}C'$ de $\operatorname{Fl}M'$
% archive sheet number: 86
\scanpage{88}
\mspage{75}

par
\[
\left\{\begin{aligned} W'&=\unsure{u|v}^{*-1}(W)\\ C'&=\text{clôture stable par coch. de base et composition}\\ &\hspace{1.5em}\text{transfinie de }u_!(C)
\end{aligned}\right.\tag*{(4.5.1)}
\]
\deleted{\unsure{donc}}
\[
u_!(C)\subset C'\subset\widetilde{u_!(C)}=\widetilde C'\ \deleted{\unsure{\text{ensemble}}}\tag*{\unsure{(4.5.1)|(4.5.1$'$)}}
\]
Je dis que \deleted{\unsure{$C$}} $(W',C')$, et au choix $(W',\widetilde C')$, est une structure de cofibration spéciale sur $M'$. On va vérifier les différentes conditions de 4.4. La condition a) est déjà acquise.

\margin{\uline{NB} Il vaut mieux commencer par \unsure{prouver|prendre} c)}
% slanted
b). La saturation \unsure{faible|forte} de $W'$ résulte de celle de $W$, via (4.5.1). La stabilité de $C'$ par coch. de base et par composition $\unsure{(=|C=}$ transfinie, est acquise par définition de $C'$, et idem pour $\widetilde C'$ (étant Q-close à gauche). Stabilité de $\mathrm{TC}'$\inserted{$\overset{\mathrm{déf}}{=}C'\cap W'$} par cochgt. de base ? Il suffit de voir que si on a un carré cocartésien \deleted{\unsure{de}}
\diagram{\begin{tikzcd}[ampersand replacement=\&] X'_0 \arrow[r,"i'"] \arrow[d] \arrow[dr,phantom,"\text{cocart}"] \& X'_1 \arrow[d,dash] \\ X'_2 \arrow[r,dash,"j'"] \& X'_3
\end{tikzcd}}
dans $M'$, avec $i'$ dans $\mathrm{TC}'=C'\cap W'$, alors $j'\in W'$ (car on sait déjà $j'\in C'$), i.e \margin{$\Vert$}$u^*(j')\in W\unsure{^{\prime}|{}}$. Nous allons supposer que \uline{$u^*$ commute aux sommes amalgamées}, et on est donc ramené à prouver que $u^*(i')\in\mathrm{TC}$. Donc
% archive sheet number: 87
\scanpage{89}
\mspage{76}

on doit prouver le

\margin{$\vert$}\uline{Lemme} 4.5.2. \deleted{\unsure{Soit $i$ une flèche de $M'$,}} \inserted{Sous les conditions (4.5.3) et (4.5.4)} ci-dessous, on a
\[
\unsure{\psi|u|v}^*(C')\subset C,\qquad u^*(\mathrm{TC}')\subset\mathrm{TC}
\]
\margin{$\vert$}et si $C=\widetilde C$, on a $u^*(\widetilde C')\subset C$, $u^*(\mathrm{T}\widetilde C')\subset\mathrm{TC}$.

\deleted{\unsure{Comme} \unsure{Je|Il}}

Pour $u^*(C')\subset C$, il faut d'abord prouver
\[
\deleted{u^*u_!(C))\subset C}
\]
\[
u_{\deleted{\unsure{*}}}^{*}(u_!(C))\subset C,\tag*{(4.5.3.)}
\]
et je ne vois aucun moyen -- il faut le poser comme hypothèse. De plus, que $u^*$ commute aux $\unsure{\varinjlim|\underline{\lim}}$ quelconques (cochgt. de base déjà \deleted{\unsure{rencontré}} admis -- à présent on a besoin des limites inductives filtrantes)
\[
u^*\ \text{commute aux }\text{\inserted{petites}}\ \varinjlim\ \text{quelc.}\tag*{(4.5.4.)}
\]
pour conclure de (4.5.3) que
\[
\begin{aligned}
u^*(C'={}&\text{clôture de }u_!(C)\text{ par coch. de base et}\\ &\text{composition transfinie})\quad\text{est également dans }C.
\end{aligned}
\]

Passons à $u^*(\mathrm{TC}')\subset\mathrm{TC}$, par ce qui précède \inserted{$u^*(C')\subset C$,} il suffit de \inserted{prouver} $u^*(\mathrm{TC}')\subset W$, mais $u^*(\mathrm{TC}')\subset u^*(W')\subset W$ par définition de $W'$.

\margin{\tikz[x=1ex,y=1ex]{\draw (0,0)--(1.15,0)--(1.15,.95); \draw (1,0)--(1,1.25)--(2,1.4);}}Supposons $C=\widetilde C$ i.e. $C$ stable par fact. \unsure{direct|directs}, et prouvons $u^*(\widetilde C')\subset\unsure{\dot C|\bar C|\widetilde C}$; $u^*(\mathrm{T}\widetilde C')\subset\unsure{\mathrm{TC}|\mathrm{T}7|\mathrm{T}\mathcal C}$. Par c) plus bas, $\widetilde C'$ est formé des fact. directs
% archive sheet number: 88
\scanpage{90}
\mspage{77}

d'éléments de $C'$, donc $\unsure{u|v}^*(C')\subset C$ implique $u^*(\widetilde C')\subset C$ puisque $C$ stable par fact. dir. Prouvons enfin $u^*(\mathrm{T}\widetilde C')\subset\mathrm{TC}$\inserted{($\overset{\mathrm{déf}}{=}C\cap W$)}. \deleted{\unsure{Comme}} $\mathrm{T}\widetilde C'=\widetilde C'\cap W'$, cela résulte de $u^*(\widetilde C')\subset C$, et de $u^*(W')\subset W$.

Prouvons b$'$). La stabilité de $C'$\inserted{ et de $\widetilde C'$,} par composition transfinie est acquise, il reste à prouver \deleted{\unsure{que}} stabilité de $\mathrm{TC}'=C'\cap W'$ et de $\mathrm{T}\widetilde C'=\widetilde C'\cap W'$. Il suffit de voir que le composé transfini \inserted{$f$} en question est dans $W'$, donc que $u^*(f)\in W$, mais cela résulte du fait que $u^*$ commute à la compos. transfinie par (4.5.4), et que $u^*$ transforme $\mathrm{TC}'$ et $\mathrm{T}\widetilde C'$ dans $\mathrm{TC}$.

\uline{Prouvons c}\inserted{ pour $C'$ (donc aussi pour $\widetilde C'$)}. Soit $C_0$ une petite partie de $C$ telle que $C\subset\widetilde C_0$, i.e $C$ est contenu dans la clôture de $C_0$ par les trois op. de \unsure{coch.} de base, de composition transfinie, de fact. directs. On voit alors que $u_!(C)$ est dans la clôture de $u_!(C_0)$ pour \unsure{les|ces} opérations, donc aussi $C'$, et on peut donc prendre $C'_0=u_!(C_0)$.

\uline{Prouvons enfin d)}, c'est la $=$ assertion pour $C'$ et pour $\widetilde C'$, puisque
\[
C'_*= \widetilde C'_* \bigl(\overset{\mathrm{déf}}{=}\mathrm{TF}'\bigr),
\]
% archive sheet number: 8[3 ou 9]
\scanpage{91}
\mspage{78}

il faut donc prouver
\[
\mathrm{TF}'\overset{\mathrm{déf}}{=}C'_*\subset W'.
\]
Cela va résulter du

\uline{Lemme} 4.5.3. On a $\mathrm{TF}'=\unsure{u|v}^{\unsure{*|\deleted{*}}{}^{-1}}(\mathrm{TF})$

C'est formel par définition de $\mathrm{TF}'$ comme
\[
\mathrm{TF}'=\bigl(u_!(\deleted{\unsure{\mathrm T}}C)\bigr)_*,
\]
de $\mathrm{TF}$ comme
\[
\mathrm{TF}=C_*,
\]
et de l'adjonction de $u_!,u^*$. Comme $\mathrm{TF}\subset W$, on en conclut
\[
\mathrm{TF}'=u^{*-1}(\mathrm{TF})\subset u^{*-1}(W)=W',
\]
\unsure{qed.}

Donc on a prouvé ceci :

\uline{Proposition} 4.6. Considérons un couple de foncteurs adjoints
\diagram{\begin{tikzcd}[ampersand replacement=\&,column sep=huge] M \arrow[r,shift left,"u_!"] \& M' \arrow[l,shift left,"u^*"]
\end{tikzcd}\,,}
où $M$ est muni d'une structure de cofibration \uline{spéciale} $(W,C)$. On suppose de plus \unsure{\deleted{que}|{}}les conditions suivantes satisfaites
\[
\left\{\begin{aligned} 1^\circ)&\ M'\ \text{stable par petites }\varinjlim,\ \text{ses éléments sont}\\ &\qquad\text{accessibles}\\ 2^\circ)&\ u^*\ \text{commute aux petites }\varinjlim\\ 3^\circ)&\ u^*(u_!(C))\subset C.
\end{aligned}\right.\tag*{(4.6.1)}
\]
% archive sheet number: 90
\scanpage{92}
\mspage{79}

Considérons les parties $W',C'$ de $\operatorname{Fl}M'$ définies par
\[
\left\{\begin{aligned} W'&=\unsure{u|v}^{*-1}(W)\\ C'&=\text{clôture de }u_!(C)\ \text{par cochgt.}\\ &\quad\text{de base et par composition}\\ &\quad\text{transfinie}
\end{aligned}\right.\tag*{(4.6.2)}
\]
Alors $(M',W',C')$ est une structure à cofibrations spéciale. \deleted{Si} Si de plus $C=\widetilde C$ $C^{\deleted{\unsure{\prime}}}$ est stable par facteurs directs, \margin{\uline{NB} C'est là un fait général pour les cat. à cofibrations spéciales, et \unsure{devrait} s'appliquer plus bas\inserted{ (\unsure{add.|cf.} \unsure{5.7.|3.7.3})}. L'hypothèse sur $C$ est inutile ici. [ou \unsure{plutôt|cf. plus bas}, il suffit \unsure{au lieu|en termes} \unsure{de cela|de 4.7} que $W$ soit stable \unsure{par} facteurs directs \uline{et} par composition transfinie, ce qui \unsure{n'est pas|est} \unsure{pour} une hypothèse plus \unsure{rigoureuse|vigoureuse} \unsure{si on y tient|si on}\dots) :}
% slanted
\deleted{\unsure{donc $C'$}} alors $(M',W',\widetilde C')$ est égal\textsuperscript{mt} une structure \deleted{\unsure{de}} à cofibrations spéciale.

\deleted{\uline{\unsure{Proposition|Remarque} \illegible{1}}}

J'aimerais, \unsure{\deleted{en}|en} vue d'énoncés du type 3.7 (établissant l'existence de foncteurs du type $\widetilde u_!$), avoir une prise sur $\mathrm{TC}'=C'\cap W'$, et sur $\mathrm{T}\widetilde C'=\widetilde C'\cap W'$, de façon précise, j'aimerais avoir la relation
\[
u_!(\mathrm{TC})\subset\mathrm{TC}'\subset u_!(\mathrm{TC})^{\sim}\tag*{(4.7.1.)}
\]
\unsure{voire} la relation légèrement plus forte
\[
\underbrace{u_!(\mathrm{TC})\unsure{\in|\subset}\mathrm{TC}'}\ \overset{\text{\unsure{tautol.|trivial.}}}{\subset}\ \underbrace{\mathrm{T}\widetilde C'\subset\bigl(u_!(\mathrm{TC})\bigr)^{\sim}},\tag*{(4.7.2)}
\]
qui se \inserted{dé}compose \deleted{\unsure{en}} en deux
\[
\left\{\begin{aligned} u_!(\mathrm{TC})&\subset\mathrm{TC}'\\ \mathrm{T}\widetilde C'&\subset\bigl(u_!(\mathrm{TC})\bigr)^{\sim}.
\end{aligned}\right.\tag*{(\unsure{4|3}.7.3)}
\]
La \deleted{plus} moins délicate des deux est la
% archive sheet number: 91
\scanpage{93}
\mspage{80}

première $\unsure{u|v}_!(\mathrm{TC})\subset\mathrm{TC}'$, qui revient, en vertu des définitions, \unsure{du fait|et du fait} $u_!(\mathrm{TC})\subset u_!(C)\subset C'$, à
\[
u_!(\mathrm{TC})\subset W'
\]
i.e par définition de $W'$, à
\[
u^*u_!(\mathrm{TC})\subset\mathrm{TC},\tag*{(4.\unsure{7|2}.4)}
\]
ce qui \deleted{\unsure{est}} est une hypothèse du type (4.6.1.) $3^\circ$), à joindre aux hypothèses (4.6.1).

La relation hypothétique
\[
\begin{aligned}
\underset{\substack{\Vert\\\widetilde C'\cap W'}}{\mathrm{T}\widetilde C'}&\subset\bigl(u_!(\mathrm{TC})\bigr)^{\sim},\qquad\text{ou seulement}\\ &\hspace{6em}\mathrm{TC}'\subset\bigl(u_!(\mathrm{TC})\bigr)^{\sim}
\end{aligned}\tag*{(a)}
\]
m'apparaît de nature nettement plus délicate, \deleted{\unsure{si|or}} \unsure{comme|l'on a} $\mathrm{T}\widetilde C'\supset\mathrm{TC}'\supset u_!(\mathrm{TC})$, elle équivaut à l'\uline{égalité} \deleted{\unsure{de} \illegible{1}} de
\[
(\mathrm{T}\widetilde C')_*\text{\inserted{, resp. $(\mathrm{TC}')_*$,}}\quad\text{et}\quad\bigl(u_!(\mathrm{TC})\bigr)_*=u^{*{}^{-1}}\bigl((\mathrm{TC})_*\bigr).
\]

Posons
\[
\left\{\begin{aligned} (\mathrm{TC})_*&\overset{\mathrm{déf}}{=}F\\ (\mathrm{TC}')_*&\overset{\mathrm{déf}}{=}F'\ \deleted{\unsure{\subset}}\supset\bigl(\mathrm{T}\widetilde C'\bigr)_*=F'_{\sim}
\end{aligned}\right.\tag*{(4.7.5)}
\]
On a donc, en vertu de l'inclusion
\[
u_!(\mathrm{TC})\subset\mathrm{TC}'\ (\subset\mathrm{T}\widetilde C')
\]
les relations
\[
F'_{\sim}\subset F'\subset u^{*-1}(F),\tag*{(4.7.6)}
\]
et la question est donc si on a égalité $F'=u^{*-1}(F)$, voire $F'_{\sim}=u^{*-1}(F)$ (d'où $F'_{\sim}=F'$).
% archive sheet number: 92
\scanpage{94}
\mspage{81}

En fait, on montre assez facilement qu'on a $F'_{\sim}=F'$, i.e. $\mathrm{T}\widetilde C'\overset{\mathrm{déf}}{=}\widetilde C'\cap W'\subset(C'\cap W'\unsure{]|)}^{\sim}$ -- c'est une histoire générale de cat. à cofibra. spéciale.\textsuperscript{*} \margin{cf prop. 5.1 ci-dessous}
% slanted
Mais je ne vois guère comment prouver (sous des hyp. convenables) que $F'=\unsure{u|v}^{*-1}(F)$. \margin{voir cependant prop. 5.8. ci-dessous}
% slanted
Il me faudra revenir là-dessus au besoin. Pour l'instant, je me bornerai à récolter le fruit de 4.6. dans un cas particulier :

\uline{Proposition} 4.8. Soit $(M,W,C)$ une catégorie à \unsure{cofibrations|cofibration} spéciale (déf. 4.4. page 74). Soit $I$ objet de $\mathrm{Cat}$, $I_0$ la cat. discrète associée, considérons $u:I_0\unsure{\mathbin{-}|\longrightarrow}I$ et
\diagram{\begin{tikzcd}[ampersand replacement=\&,column sep=huge] M(I_0) \arrow[r,shift left,"u_!"] \& M(I) \arrow[l,shift left,"u^*"]
\end{tikzcd}\,,}
enfin munissons $M(I_0)\simeq M^{I_0}$ de la structure de \inserted{cat. à} cofibration, \uline{produit} de celles des facteurs, i.e. définissons $W(I_0)$, $C(I_0)$ \deleted{\unsure{pour $F_0$,|pour $E_0$,}} ``fibre à fibre''. (Cela coïncide avec la définition générale donnée au \S\,2, pour $I_0$ muni de la structure d'ordres discrète, et fait de $M'(I_0)$ une \deleted{structure} cat. à cofibration spéciale.) Alors le couple $(u_!,u^*)$ \unsure{satisfais|satisfait}
% archive sheet number: 93
\scanpage{95}
\mspage{82}

aux conditions de la prop. 4.6, \deleted{\unsure{et}} ainsi qu'à la condition supplémentaire à (4.6.1.)
\[
4^\circ)\quad\unsure{u|v}^*(u_!(\mathrm{TC}(\unsure{I|1}))\subset\mathrm{TC}(I).\tag*{(4.8.1)}
\]
Donc on trouve une structure de cat. à cofibr. spéciale sur $M(I)$, en posant
\[
W(I)=u^{*-1}(W(I_0))\tag*{(4.8.2)}
\]
(c'est en fait \inserted{\unsure{équiv. à|égal à}} la définition habituelle de $W(I)$) et
\[
\begin{aligned}
C(I)={}&\text{clôture }\deleted{\text{\unsure{pour}}}\ u_!(C(I_0))\ \text{pour}\\ &\text{les op. de coch. de base et de}\\ &\text{composition transfinie.}
\end{aligned}\tag*{(4.8.3)}
\]

On a \inserted{de plus,} sous ces conditions :
\[
\underset{\substack{\Vert\,\mathrm{déf}\\(C(I))_*}}{\mathrm{TF}(I)} =(u^*)^{-1}\bigl(\underset{\substack{\Vert\,\mathrm{déf}\\(C_{\inserted{I_0}}\deleted{\unsure{(F)}})_*,\ \text{et}\\\text{\unsure{définit|défini} \unsure{aussi|ainsi}}\\\text{fibre par fibre,}}}{\underline{\mathrm{TF}(I_0)}}\bigr)\tag*{(4.8.8)}
\]
\margin{Ces inclusions devraient être des égalités, mais je ne sais pas \unsure{le|les} prouver sans faire des hypothèses plus fortes, du type \unsure{\S\,9|p.~9}, \uline{cat. de Quillen spéciales}.}
% slanted
\[
\left\{\begin{aligned} u_!\bigl(\overbrace{\mathrm{TC}(I_0)}^{C(I_0)\cap W(I_0)}\bigr)&\subset\mathrm{TC}(I)\\ F(I)\overset{\mathrm{déf}}{=}\bigl(\mathrm{TC}(I)\bigr)_*&\subset(u^*)^{-1}\bigl(F(I_0)\bigr).
\end{aligned}\right.\tag*{(4.8.9)}
\]

\uline{\unsure{Dém|Preuve}} On sait déjà que $M(I)$ est stable par \inserted{petites} $\varinjlim$ \deleted{$M$ l'étant,} et que \deleted{\unsure{tout|les}} \unsure{élément|l'élément} est accessible, \unsure{comme} il résulte aussitôt des propriétés idoines de $M$. Il est clair aussi que $u^*$ commute aux \inserted{petites} $\varinjlim$. Il reste à prouver (4.6\unsure{.}1. $3^\circ$), et $4^\circ$ de (4.8.1) \unsure{\deleted{q}|{}}\unsure{i.e.}
% archive sheet number: 94
\scanpage{96}
\mspage{83}

que $\unsure{u|v}_{\unsure{\cdot|{}}}^*u_!:M(I_0)\longrightarrow M(I_0)$ \deleted{\unsure{transforme}} \inserted{applique} $C$ dans $C$, $\mathrm{TC}$ dans $\mathrm{TC}$. Or on a
\[
u_!(F_0)(i)=\varinjlim_{(\alpha,\deleted{\unsure{i}}j)\in i\backslash I_0}F_0(j)
\]
\margin{$\vert$ où \deleted{\unsure{$\alpha:u(i)\longrightarrow j=j$}} \inserted{$j\in\mathrm{Ob}\,I_0=\deleted{\unsure{\mathrm{Ob}}}$} $\mathrm{Ob}\,I$, et \deleted{$\alpha:i\longrightarrow$} $\alpha:i\longrightarrow u(j)=j$} et $i\backslash I_0$ est la catégorie discrète \inserted{au dessus de $I_0$}, dont la fibre en $j\in\mathrm{Ob}\,I_0=\mathrm{Ob}\,I$ est $\mathrm{Hom}_I(i,j)$ :
\[
i\backslash I_0=\coprod_{j\in\mathrm{Ob}\,I}\binom{\text{cat. discrète associée}}{\text{à }\mathrm{Hom}_I(i,j)}\,,
\]
donc
\diagram{\begin{tikzcd}[ampersand replacement=\&,column sep=small,row sep=small] u_!(F_0)(i)={} \& \displaystyle\coprod_{j\in\underset{\substack{\Vert\\\mathrm{Ob}\,I_0}}{\mathrm{Ob}\,I}} \underbrace{F_0(j)^{(\mathrm{Hom}_I(i,j))}}_{\substack{\text{somme directe}\\\text{de }\mathrm{Hom}_I(i,j)\\\text{copies de }F_0(j)}} ,\\ \text{\deleted{\unsure{donc}} on trouve} \& {}\\ \text{et}\quad u^*u_!(F_0)(\underset{\substack{\rotatebox[origin=c]{-90}{\(\in\)}\\\mathrm{Ob}\,I_0=\mathrm{Ob}\,I}}{i})=u_!(F_0)(i)={} \& {} \arrow[uu,bend right=23]
\end{tikzcd}\,,}

Comme $C\deleted{\unsure{(I_0)}}$ et $\mathrm{TC}\deleted{\unsure{(I_0)}}$ sont stables par sommes directes quelconques (l'étant par sommes directes finies, à cause des conditions de stabilité \unsure{finie|finies} sur $C$ et $\mathrm{TC}$, et puis par sommes infinies à cause de l'axiome \margin{\unsure{mérite} un lemme général \ldots}
% slanted
de stabilité transfinie \ldots), il s'ensuit \inserted{bien} que $u^*u_!$ \deleted{\unsure{transforme}} \inserted{applique} $C(I_0)$ et $\mathrm{TC}(I_0)$ en eux-mêmes,

Ainsi, la \deleted{\unsure{conclusion}} première assertion de 4.8. résulte de 4.6. De plus (4.\unsure{6|8}.8.) et (4.8.\unsure{4|9}) ont été vus dans 4.7., dans les conditions générales de 4.6 : et en supposant de plus le $4^\circ$ de (4.7.4)\unsure{)|{}}.
% archive sheet number: 95
\scanpage{97}
\mspage{84}

5. Relations des cat. à cofibrations spéciales avec les \deleted{\unsure{Structures|structures}} catégories de Quillen spéciales.

Soit $(M,W,C)$ une catégorie à cofibrations spéciale. J'introduis
\[
\mathrm{TF}=C_*,\quad \mathrm{TC}=C\cap W,\quad F=(\mathrm{TC})_*,\tag*{(5.1)}
\]
donc on a
\diagram{$(5.2)$\qquad\begin{tikzcd}[ampersand replacement=\&] \mathrm{TC} \arrow[r,hook] \arrow[d,<->] \& C \arrow[d,<->,"{-}"{description,fill=none,inner sep=0pt}]\\ F \& \mathrm{TF} \arrow[l,hook']
\end{tikzcd},}
$(C,\mathrm{TF})$ étant un \unsure{précouple} de Quillen clos à droite.

\deleted{\uline{Proposition \unsure{5.1}}} On sait alors \inserted{(\uline{$\overline{\mathrm{XV}}$} 4.\unsure{2|7}. cor. 1 \unsure{$\sim$|v.|bis} page 31)} que $\widetilde C$ est formé des facteurs directs d'éléments de $C$. Posons
\[
\mathrm{T}\widetilde C=\widetilde C\cap W\tag*{(5.3)}
\]
d'où
\[
\mathrm{TC}\subset\mathrm{T}\widetilde C\tag*{(5.4)}
\]

\uline{Proposition} 5.1. \inserted{Soit $(M,W,C)$ une cat. à cofibrations spéciale.} On a $\mathrm{T}\widetilde C\subset(\mathrm{TC})^{\sim}$\inserted{ (et \unsure{$=$} $\mathrm{T}\widetilde C\subset(\mathrm{TC})^{\unsure{\leftrightarrows|\leftrightarrow}}$)\textsuperscript{*}}, \unsure{ou} ce qui revient au $=$ en vertu de (5.4), \deleted{\unsure{$\mathrm{TC}$}} $\mathrm{TC}$ et $\mathrm{T}\widetilde C$ ont $=$ Q-clôture à gauche : \margin{\textsuperscript{*}\uline{NB} Pour une partie $A$ de $\operatorname{Fl}M$, $A^{\leftrightarrows}$ désigne l'ens des facteurs directs d'éléments de $A$.}
% slanted
\[
\mathrm{T}\widetilde C\subset(\mathrm{T}\widetilde C)^{\sim}=(\mathrm{TC})^{\sim}\ \text{\inserted{(et $=$ $\mathrm{T}\widetilde C\subset(\mathrm{TC})^{\leftrightarrows}$)}}\,\text{d'où}\tag*{(5.5.)}
\]
\inserted{
\[
\boxed{(\widetilde C)_*=C_*(=\mathrm{TF}),\quad(\mathrm{T}\widetilde C)_*=(\mathrm{TC})_*\ (=F).}\tag*{(5.6)}
\]
}
\margin{Ces relations ne dépendent pas de l'hyp. $W$ stable par fact. directs.}
% slanted
De plus, \inserted{si $W$ stable par facteurs directs (ou si seulement les \unsure{fact. dir. d'éléments} de $\mathrm{TC}$ sont $\in W$),} \deleted{\unsure{$(M,W,C)$ est une catégorie}} \inserted{et si de plus $W$ est stable par composition transfinie,} alors $(M,W,\widetilde C)$ \deleted{\illegible{1}} est une cat. à cofibrations spéciale,

\deleted{Enfin, si $W$ est stable par \unsure{composition} \unsure{et par fact. directs} \inserted{\unsure{comp$^{\mathrm{n}}$ de facteurs directs et de flèches de $\mathrm{TC}$ est dans $W$}}} \uline{De plus} $\mathrm{T}\widetilde C$ est Q-\deleted{\illegible{1}}\inserted{\unsure{saturée|saturé}} à gauche, \unsure{donc} (par 5.5)
\[
\mathrm{T}\widetilde C=(\mathrm{TC})^{\leftrightarrows}.\tag*{\deleted{\unsure{(5.6)}}(5.7)}
\]
% archive sheet number: 96
\scanpage{98}
\mspage{85}

\uline{Dém} Prouvons
\[
\mathrm{T}\widetilde C\subset(\mathrm{TC})^{\sim}
\]
et plus précisément, que tout élément de $\mathrm{T}\widetilde C$ est facteur direct d'un élément de $\mathrm{TC}$ :
\[
\mathrm{T}\widetilde C\subset(\mathrm{TC})^{\unsure{\leftrightarrows|\leftrightarrow}}.
\]
Soit donc $f\in\mathrm{T}\widetilde C\overset{\mathrm{déf}}{=}\widetilde C\cap W$, considérons une
\diagram{\begin{tikzcd}[ampersand replacement=\&] X \arrow[rr,"f\in C\cap W"] \arrow[dr,"i\in C"'] \& {} \& Y \arrow[dl,dashed,bend right=16,"g"']\\ {} \& X' \arrow[ur,"p\in\mathrm{TF}"'] \& {}
\end{tikzcd}}
factorisation de $f$ en $pi$, avec $i\in C$, $p\in\deleted{\unsure{W}}\mathrm{TF}\deleted{\unsure{\subset}}$. Comme $\mathrm{TF}\subset W$, on a $p\in W$, donc $i\in W$, donc $i\in C\cap W\overset{\mathrm{déf}}{=}\mathrm{TC}$. Comme $f\in C$, $p\in\mathrm{TF}$, il existe $g:Y\unsure{\mathbin{-}|\to}X'$ tel que $pg=\mathrm{id}_Y$, $gf=i$, donc $f$ apparaît comme facteur direct de $i\in\mathrm{TC}\deleted{\unsure{\cap W}}$, qed.

Prouvons que $(M,W,\widetilde C)$ est une cat. à cofibration spéciale (déf. 4.4. page 74). Condition a) sur $M$ acquise, condition b), \inserted{la partie} sur $W$ aussi. La stabilité de $\widetilde C$ par cochgt de base et et par composition OK, aussi (pour b$'$) la stabilité par comp. transfinie. Stabilité de \inserted{$\mathrm{T}\widetilde C=\widetilde C\cap W$} par cochgt de base : il reste à prouver que si on a carré cocartésien
\diagram{\begin{tikzcd}[ampersand replacement=\&] X \arrow[r,"{\unsure{f|i}}"] \arrow[d,dash] \arrow[dr,phantom,"\text{\unsure{cocart.}}"] \& Y \arrow[d,dash,"{\unsure{h|{}}}"]\\ X' \arrow[r,"{\deleted{\unsure{f}}i'}"'] \& Y'
\end{tikzcd},}
si $i\in\widetilde C\cap W$, alors $i'\in\deleted{\unsure{C}} W$ (car on sait déjà que $i'\in\widetilde C$). Mais la démonstration précédente montre que $i$ est facteur direct \inserted{\uline{dans} \uline{$X\backslash M$}} d'une flèche de $\mathrm{TC}$, donc $i'$ est facteur direct d'une flèche de $\mathrm{TC}\subset W$. \deleted{\unsure{a fortiori une flèche|$i'$ factorise une flèche}} Donc si $W$ est stable par facteurs directs, \deleted{\unsure{alors}} $i'\in W$.
% archive sheet number: 97
\scanpage{99}
\mspage{86}

\deleted{La première des relations (5.6) est une tautologie, la deuxième \deleted{pour} équivaut au fait que $\mathrm{TC}$ et $\mathrm T\widetilde C$ ont $\mathord{=}$ Q-clôture à gauche, déjà prouvé (5.5.).}

Cela achève de prouver b). Pour b'), il reste à voir que $\mathrm T\widetilde C$ est stable par composition transfinie. Je ne vois pas d'autre façon de le voir, qu'en supposant que $W$ lui-$=$ est stable par composition transfinie (ce qui paraît à peine moins exigeant que la stabilité de $W$ par $\varinjlim$). \inserted{Cela achève de prouver b,b').} La condition c) sur $C$ \deleted{\unsure{signi|sign}} implique la $=$ condition sur $\widetilde C$. Enfin, comme $\widetilde C_{\unsure{*|\alpha}}=C_*$, la condition d) est la $=$ pour $C$ et $\widetilde C$. Cela achève de prouver que $(\unsure{\mathcal M|M},W,\widetilde C)$ est une cat. à cofibrations spéciale. (\uline{NB} On avait besoin des \deleted{\unsure{2 hyp|2 hy}} deux hyp. suppl. sur $W$, stabilité par facteur direct et par composition transfinie, pour établir que $\mathrm T\widetilde C$ est stable par cochg\textsuperscript{t} de base et par composition transfinie). La dernière assertion de 5.1 \unsure{de|et} la relation (5.7) qui en résulte, \unsure{revient} à dire que $\mathrm T\widetilde C=\widetilde C\cap W$ est stable par facteurs directs, ce qui est évident par la $=$ \deleted{\unsure{hyp.}} propriété de $\widetilde C$, et l'hyp. faite sur $W$ de \deleted{que les facteurs directs d'\unsure{élém}} stabilité par fact. directs, (ou seulement que les fact. directs d'éléments de $\mathrm T\widetilde C$ sont dans $W$, donc dans $\mathrm T\widetilde C$).
% archive sheet number: 98
\scanpage{100}
\mspage{87}

\uline{\deleted{\unsure{Corollaire}}}\inserted{Proposition} 5.2. Soit $(M,W,C)$ une catégorie à cofibrations spéciale. \deleted{Pour que \unsure{le|la}} Alors avec les notations de (5.1), on a
\[
 W=\mathrm{TF}\circ\mathrm{TC} \tag*{(5.8)}
\]
\[
 \mathrm{TF}=F\cap W. \tag*{(5.9)}
\]

\uline{Dém} Soit $f\in W$, écrivons $f=pi$, avec $i\in C$, $p\in\mathrm{TF}$, donc $p\in W$, donc $i\in W$, donc $i\in C\cap W=\mathrm{TC}$, d'où (5.8).

Prouvons (5.9). On a $\mathrm{TF}\subset F\cap W$, prouvons inversement qu'une $f\in F\cap W$ est dans $\mathrm{TF}$.
\diagram{\begin{tikzcd}[ampersand replacement=\&]  X \arrow[rr,"f\in F\cap W"] \arrow[dr,"i\in\mathrm{TC}"'] \& \& Y\\  \& X'\arrow[ur,"p\in\mathrm{TF}"'] \&
\end{tikzcd}}
Écrivons $f=pi$, avec $p\in\mathrm{TF}$, $i\in\mathrm{TC}$. Comme $f\in F$, $\deleted{\unsure{c}}i\in\mathrm{TC}$, et $\mathrm{TC}\longleftrightarrow F$ par définition de $F$, $\exists\,r:X'\unsure{\mathbin{-}|\to}X$ tel que $ri=\mathrm{id}_X$, $fr=p$. Donc $f$ est facteur direct de $p\inserted{\in\mathrm{TF}}$, et comme $\mathrm{TF}=C_*$  est stable par facteurs directs, on conclut que $f\in\mathrm{TF}$.

\uline{Corollaire} 5.3. Soit $(M,W,C)$ une catégorie à cofibrations spéciale. Pour que $(W,C,F\overset{\mathrm{déf}}{=}(\mathrm{TC})_*)$ soit un triple de Quillen dans $M$, il f. et il suffit que \deleted{$\unsure{\mathrm{TC},F|\mathrm{TC},E}$} le couple $(\mathrm{TC},F)$ satisfasse à l'axiome de factorisation, (i.e. que ce soit un précouple de Quillen). Pour que cette structure de Quillen soit \unsure{clos|close}, il f. et il s. que $C$ \deleted{et $W$} \unsure{soit|soient} stable par par facteurs directs.
% archive sheet number: 99
